\SourcePage{654}{657}
\section{Scattering of slow particles}

Consider elastic scattering in the limiting case of low particle velocities.
The wavelength is assumed large compared with the range $a$ of $U(r)$
($ka\ll1$), and the energy small compared with the field within that range.
We must determine the limiting dependence of the phases $\delta_l$ on small
$k$.

\SourcePage{655}{658}
For $r\lesssim a$, only the $k^2$ term may be neglected in the exact
Schrödinger equation (123.7):
\begin{equation}
 R_l''+\frac2rR_l'-\frac{l(l+1)}{r^2}R_l
 =\frac{2m}{\hbar^2}U(r)R_l. \tag{132.1}
\end{equation}
For $a\ll r\ll1/k$, the $U(r)$ term may also be omitted:
\begin{equation}
 R_l''+\frac2rR_l'-\frac{l(l+1)}{r^2}R_l=0. \tag{132.2}
\end{equation}
Its general solution is
\begin{equation}
 R_l=c_1r^l+\frac{c_2}{r^{l+1}}. \tag{132.3}
\end{equation}
The constants $c_1,c_2$, different for each $l$, can in principle be found
only by solving (132.1) for the specified $U(r)$.

At $r\sim1/k$, $U(r)$ may be omitted but $k^2$ may not:
\begin{equation}
 R_l''+\frac2rR_l'+\left[k^2-\frac{l(l+1)}{r^2}\right]R_l=0. \tag{132.4}
\end{equation}
The free-motion solution is
\begin{equation}
 R_l=c_1\frac{(-1)^l(2l+1)!!}{k^{2l+1}}r^l
 \left(\frac1r\frac d{dr}\right)^l\frac{\sin kr}{r}
 +c_2\frac{k^l}{(2l-1)!!}
 \left(\frac1r\frac d{dr}\right)^l\frac{\cos kr}{r}. \tag{132.5}
\end{equation}
The coefficients make this reduce to (132.3) for $r\ll1/k$, matching the
solutions in the regions $kr\ll1$ and $kr\sim1$.

For $kr\gg1$ it becomes
\[
 R_l\sim\frac{c_1(2l+1)!!}{rk^{l+1}}\sin\left(kr-\frac{\pi l}2\right)
 +\frac{c_2k^l}{r(2l-1)!!}\cos\left(kr-\frac{\pi l}2\right),
\]
or
\begin{equation}
 R_l\sim\mathrm{const}\cdot\frac1r
 \sin\left(kr-\frac{\pi l}2+\delta_l\right). \tag{132.6}
\end{equation}

\SourcePage{656}{659}
Here
\begin{equation}
 \tan\delta_l=\frac{c_2}{c_1(2l-1)!!(2l+1)!!}k^{2l+1}; \tag{132.7}
\end{equation}
all phases are small at small $k$. By (123.15),
$f_l=(e^{2i\delta_l}-1)/(2ik)\approx\delta_l/k$, and hence
\begin{equation}
 f_l\sim k^{2l}. \tag{132.8}
\end{equation}
All $l\ne0$ amplitudes are therefore small compared with the $l=0$ or
$s$-wave amplitude. Neglecting them,
\begin{equation}
 f(\theta)\approx f_0=\frac{\delta_0}{k}=\frac{c_2}{c_1}=-a, \tag{132.9}
\end{equation}
so $d\sigma=a^2do$ and
\begin{equation}
 \sigma=4\pi a^2. \tag{132.10}
\end{equation}
Slow-particle scattering is isotropic and energy-independent\footnote{For
electron scattering by atoms, the length compared with $1/k$ is the atomic
radius, which in complex atoms reaches several Bohr radii. Its large value
restricts constant cross-sections to energies of fractions of an electron
volt. At higher energies the cross-section depends strongly on energy (the
Ramsauer effect).}. The constant $a$, positive or negative, is the
\emph{scattering length}.

These arguments assume sufficiently rapid decrease of $U(r)$. At large $r$,
the second term of (132.3) is small. Retaining it is legitimate only if the
terms $\sim c_2/r^{l+1}r^2$ kept in (132.2) exceed the omitted term
$UR_l\sim Uc_1r^l$. Thus $U(r)$ must

\SourcePage{657}{660}
decrease faster than $1/r^{2l+3}$ for (132.8) to hold. In particular, (132.9)
and energy-independent isotropic scattering require decrease faster than
$1/r^3$.

If $U$ decreases exponentially, further terms in the power expansion of
$f_l$ can be characterized. Section 128 showed that $f_l(E)$ is real for real
negative $E$\footnote{At small $E$, (128.6) already holds for decrease as
$e^{-r/a}$.}. The same is true of $g_l(E)$ in $f_l=1/(g_l-ik)$, since $ik$ is
real for $E<0$; by definition $g_l$ is also real for $E>0$. Hence $g_l$ has an
expansion in integral powers of $E$, or even powers of $k$. Consequently
$f_l(k)$ expands in integral powers of $ik$: even powers have real
coefficients and odd powers imaginary ones. By (132.8), $f_l$ begins with
$k^{2l}$, and $g_l$ begins with $k^{-2l}$.

For a power-law tail $U\sim\beta r^{-n}$ with $n\leq3$, the constant-amplitude
result fails. For $n<1$, at sufficiently low velocity and practically every
impact parameter $\rho$,
\begin{equation}
 \rho|U(\rho)|\gtrsim\hbar v, \tag{132.11}
\end{equation}
so scattering is classical. For $1<n<2$, this holds over a substantial range
of moderate $\rho$, giving classical scattering through appreciable angles,
while another range satisfies

\SourcePage{658}{661}
\begin{equation}
 \rho|U(\rho)|\ll\hbar v, \tag{132.12}
\end{equation}
the perturbative condition. For $n>2$, at large distances
\begin{equation}
 |U|\ll\frac{\hbar^2}{mr^2}, \tag{132.13}
\end{equation}
so the distant contribution is perturbative even if the nearer region is
not\footnote{Low-velocity scattering never becomes semiclassical here,
because (132.11) is incompatible with the simultaneous condition
$|U(\rho)|\gtrsim E$.}. Choose $r_0\ll1/k$ so that (132.13) holds for
$r\gg r_0$. By (126.12), the contribution from this region is
\begin{equation}
 -\frac{2m\beta}{\hbar^2q}\int_{r_0}^{\infty}\frac{\sin qr}{r^{n-1}}dr
 =-\frac{2m\beta}{\hbar^2}q^{n-3}
 \int_{qr_0}^{\infty}\frac{\sin x}{x^{n-1}}dx. \tag{132.14}
\end{equation}
For $2<n<3$ the lower limit is convergent and may be replaced by zero at
$kr_0\ll1$, making this the dominant contribution:
\begin{equation}
 f\sim q^{-(3-n)},\qquad2<n<3. \tag{132.15}
\end{equation}
For $n=3$ the lower-limit divergence is logarithmic and again dominant:
\begin{equation}
 f\sim\ln\frac{\mathrm{const}}q,\qquad n=3. \tag{132.16}
\end{equation}
For $n>3$ the distant contribution vanishes as $k\to0$ and (132.9) governs.
Nevertheless (132.14) is anomalous: repeated integration by parts separates
even powers of $k$, leaving a convergent term proportional to $k^{n-3}$,
which is not generally even.

\SourcePage{659}{662}
\footnote{If $n=2p+1$ is an odd integer, $n-3=2p-2$ is even, but (132.14)
still has an anomalous part proportional to $q^{2p-2}\ln q$.}

\subsection*{Problems}

\ProblemHeading{1.} Find the slow-particle cross-section for a spherical square
well of depth $U_0$ and radius $a$.

\SolutionHeading Assume $ka\ll1$ and $k\ll\varkappa$, where
$\varkappa=\sqrt{2mU_0}/\hbar$. Only $\delta_0$ matters. With $l=0$, (132.1)
for $\chi=rR_0$ gives $\chi''+\varkappa^2\chi=0$ inside, whose regular solution
is $\chi=A\sin\varkappa r$. Outside,
$\chi=B\sin(kr+\delta_0)$. Continuity of $\chi'/\chi$ at $a$ gives
\[
 \varkappa\cot a\varkappa=k\cot(ka+\delta_0)
 \approx\frac{k}{ka+\delta_0},
\]
and hence\footnote{This formula fails when $\varkappa a$ is close to an odd
multiple of $\pi/2$. Then the discrete spectrum has a negative-energy level
near zero (see Problem 1 of \S~33), and the formulas of the next section are
required.}
\[
 f=\frac{\tan\varkappa a-\varkappa a}{\varkappa}.
\]
For $\varkappa a\ll1$, it gives
$\sigma=(4/9)\pi a^2(\varkappa a)^4$, agreeing with Born approximation.

\ProblemHeading{2.} The same for a spherical square potential barrier of
height $U_0$.

\SolutionHeading Replace $U_0$ by $-U_0$, and hence $\varkappa$ by
$i\varkappa$:
\[
 f=\frac{\tanh\varkappa a-\varkappa a}{\varkappa}.
\]

\SourcePage{660}{663}
For $\varkappa a\gg1$, $f=-a$ and $\sigma=4\pi a^2$. This is scattering by an
impenetrable sphere; classical mechanics would give one quarter of this value,
$\pi a^2$.

\ProblemHeading{3.} Find the low-energy cross-section in
$U=\alpha/r^n$, where $\alpha>0$, $n>3$.

\SolutionHeading For $l=0$, (132.1) is
\[
 \chi''-\gamma^2\frac{\chi}{r^n}=0,\qquad
 \gamma^2=\frac{2m\alpha}{\hbar^2}.
\]
With
\[
 r=\left(\frac{2\gamma}{(n-2)x}\right)^{2/(n-2)},\qquad
 \chi=\sqrt r\,\varphi,
\]
it becomes
\[
 \varphi''+\frac1x\varphi'-\left[1+\frac1{(n-2)^2x^2}\right]\varphi=0.
\]
The solution vanishing at $r=0$ is, up to a factor,
\[
 \chi=\sqrt r\,H_{1/(n-2)}^{(1)}
 \left(\frac{2i\gamma}{n-2}r^{-(n-2)/2}\right).
\]
Using the small-argument formulas for $H_p^{(1)}$, at
$\gamma\ll r\ll1/k$ one finds $\chi=\mathrm{const}(c_1r+c_2)$ and
\[
 f=-\left(\frac{\gamma}{n-2}\right)^{2/(n-2)}
 \frac{\Gamma[(n-3)/(n-2)]}{\Gamma[(n-1)/(n-2)]}.
\]

\ProblemHeading{4.} Find the slow-particle amplitude for a field with tail
$U=\beta r^{-n}$, $2<n\leq3$.

\SolutionHeading The leading term is (132.14) with lower limit zero:
\begin{equation}
 f=-\frac{\pi m\beta}{\hbar^2\Gamma(n-1)\cos(\pi n/2)}q^{n-3},
 \qquad2<n<3, \tag{1}
\end{equation}
and for $n=3$,
\begin{equation}
 f=-\frac{2m\beta}{\hbar^2}\ln\frac{\mathrm{const}}q. \tag{2}
\end{equation}
Expansion of (1) in Legendre polynomials gives
\begin{equation}
 f_l=-\frac{m\beta}{2\hbar^2}
 \frac{\Gamma[(n-1)/2]\Gamma[l-(n-3)/2]}
 {\Gamma(n/2)\Gamma[(n+1)/2+l]}k^{n-3}. \tag{3}
\end{equation}

\SourcePage{661}{664}
For $n>3$, (1) gives the anomalous part. In partial amplitudes, (3) is always
dominant when $2l>n-3$; then $f_l\sim k^{n-3}$ replaces (132.8).

\ProblemHeading{5.} Find the slow-particle amplitude in
$U(r)=-U_0e^{-r/a}$, $U_0>0$.

\SolutionHeading Put
\[
 x=2a\varkappa e^{-r/2a},\qquad \varkappa=\sqrt{2mU_0}/\hbar.
\]
Equation (132.1) for $\chi=rR_0$ becomes
$\chi_{xx}+x^{-1}\chi_x+\chi=0$, with solution
$\chi=AJ_0(x)+BN_0(x)$. The condition $\chi(0)=0$ at $r=0$ gives
\[
 A/B=-N_0(2\varkappa a)/J_0(2\varkappa a).
\]
For $a\ll r\ll1/k$, $x\ll1$ (with, of course,
$a\varkappa\exp(-1/ak)\ll1$), and
\[
 \chi\approx A+B\frac2\pi\ln\frac{\gamma x}2
 =A+B\frac2\pi\ln(a\varkappa\gamma)-\frac B{\pi a}r,
\]
where $\gamma=e^C=1{,}78\ldots$. Comparison with (132.3) gives
\[
 f=-a\left[\frac{\pi A}{B}+2\ln(a\varkappa\gamma)\right]
 =a\left[\frac{\pi N_0(2\varkappa a)}{J_0(2\varkappa a)}
 -2\ln(a\varkappa\gamma)\right].
\]
For $\varkappa a\ll1$, $f=2a^3\varkappa^2$; for
$\varkappa a\gg1$, $f=-2a\ln(\varkappa a\gamma)$.

\ProblemHeading{6.} Find the low-energy amplitude in second-order perturbation
theory (I.~Ya.~Pomeranchuk, 1948).

\SolutionHeading As $k\to0$, the integral in the second term of (130.13) is
\[
 -\int U(\mathbf k'-\mathbf k'')\frac{U(\mathbf k''-\mathbf k)}{k''^2}d^3k''
 =-2\pi^2\iint\frac{U(\mathbf r)U(\mathbf r')}{|\mathbf r-\mathbf r'|}dV\,dV',
\]
using
\[
 \int e^{i\mathbf k(\mathbf r-\mathbf r')}\frac{4\pi}{k^2}
 \frac{d^3k}{(2\pi)^3}=\frac1{|\mathbf r-\mathbf r'|}.
\]
Thus
\begin{equation}
 f=-\frac{m}{2\pi\hbar^2}\int U\,dV
 +\left(\frac{m}{2\pi\hbar^2}\right)^2
 \iint\frac{U(\mathbf r)U(\mathbf r')}{|\mathbf r-\mathbf r'|}dV\,dV'. \tag{1}
\end{equation}

\SourcePage{662}{665}
For a central field,
\[
 f=-\frac{2m}{\hbar^2}\int Ur^2dr+\frac{8m^2}{\hbar^4}
 \iint_{r'>r}U(r)U(r')r^2dr\,r'dr'.
\]
The second term in (1) is always positive. Hence for repulsion the first Born
approximation overestimates, and for attraction underestimates, the
low-energy cross-section.

\ProblemHeading{7.} Determine the energy dependence of the slow-particle
scattering amplitude in two dimensions.

\SolutionHeading At large distances the wave function is given by Problem 1
of \S~124. As in three dimensions, the $m=0$ state dominates, so $f$ is
angle-independent. Replacing $e^{ik\rho}/\sqrt\rho$ by the exact free solution
with this asymptotic form gives, for all $\rho\gg a$,
\begin{equation}
 \psi=e^{ikz}+\frac f2\sqrt{\pi k}\,H_0^{(1)}(k\rho). \tag{1}
\end{equation}
At $\rho\ll1/k$, use
$H_0^{(1)}(x)\approx(2i/\pi)\ln[2i/(\gamma x)]$ to obtain
\begin{equation}
 \psi\approx\left(1+f\sqrt{\frac{k}{\pi}}
 \ln\frac{2i}{\gamma k}\right)
 -f\sqrt{\frac{k}{\pi}}\ln\rho. \tag{2}
\end{equation}
This is the general free solution $C_1+C_2\ln\rho$ for
$1/k\gg\rho\gg a$. The real, energy-independent ratio fixed by the zero-energy
equation at $\rho\sim a$ is written
\begin{equation}
 C_1/C_2=-\ln r_0, \tag{3}
\end{equation}
where $r_0$ has dimensions of length. Comparison gives
\begin{equation}
 f=-\left[\sqrt{\frac{k}{\pi}}
 \ln\frac{2i}{\gamma kr_0}\right]^{-1}. \tag{4}
\end{equation}

\SourcePage{663}{666}
Thus, unlike the three-dimensional case, the two-dimensional cross-section
increases as the energy decreases. For scattering by an infinitely high
cylindrical barrier of radius $a$, the constant $r_0$ equals $a$.
